
\documentclass{article}
\usepackage{colortbl}
\usepackage{makecell}
\usepackage{multirow}
\usepackage{supertabular}

\begin{document}

\newcounter{utterance}

\centering \large Interaction Transcript for game `wordle', experiment `high\_frequency\_words\_no\_clue\_no\_critic', episode 1 with reflexion{-}Qwen3.5{-}27B.
\vspace{24pt}

{ \footnotesize  \setcounter{utterance}{1}
\setlength{\tabcolsep}{0pt}
\begin{supertabular}{c@{$\;$}|p{.15\linewidth}@{}p{.15\linewidth}p{.15\linewidth}p{.15\linewidth}p{.15\linewidth}p{.15\linewidth}}
    \# & \multicolumn{2}{c}{Player} && \multicolumn{2}{c}{Game Master} \\
    \hline

    \theutterance \stepcounter{utterance}  
    & & & \multicolumn{4}{p{0.6\linewidth}}{
        \cellcolor[rgb]{0.9,0.9,0.9}{
            \makecell[{{p{\linewidth}}}]{
                \texttt{\tiny{[P1$\langle$GM]}}
                \texttt{You are a language wizard who likes to guess words by using the given rules.} \\
\\ 
\texttt{Welcome to Wordle! You have six attempts to guess the target word, a valid English word of five lowercase letters (a{-}z). Please use the tags "explanation:" and "guess:" to provide a concise explanation for each guess.} \\
\\ 
\texttt{For instance, if your guess is "apple", your response should be} \\
\texttt{explanation: this is a common five{-}letter English word, and I am starting my guess with this word.} \\
\texttt{guess: apple} \\
\\ 
\texttt{After each guess, your answer will be validated, and you will receive feedback indicating which letters are correct (green), which letters are correct but in the wrong position (yellow), and which letters are incorrect (red). This feedback can be useful in determining which letters to include or exclude in your next guess.} \\
\\ 
\texttt{For example, the feedback for "apple" might be:} \\
\texttt{guess\_feedback: a<yellow> p<yellow> p<green> l<yellow> e<red>} \\
\\ 
\texttt{The explanation should contain details about how the guess\_feedback is used to arrive at a new guess.} \\
\\ 
\texttt{Do not generate any other text.} \\
\\ 
\texttt{Let's begin with your first guess.} \\
            }
        }
    }
     \\ \\

    \theutterance \stepcounter{utterance}  
    & \multicolumn{4}{p{0.6\linewidth}}{
        \cellcolor[rgb]{0.561,0.737,0.561}{
            \makecell[{{p{\linewidth}}}]{
                \texttt{\tiny{[P1$\rangle$GM]}}
                \texttt{explanation: I am starting with a common five{-}letter word containing frequent vowels and consonants to maximize the information gained from the first feedback.} \\
\texttt{guess: slate} \\
            }
        }
    }
    & & \\ \\

    \theutterance \stepcounter{utterance}  
    & & & \multicolumn{2}{p{0.3\linewidth}}{
        \cellcolor[rgb]{0.9,0.9,0.9}{
            \makecell[{{p{\linewidth}}}]{
                \texttt{\tiny{[GM$|$GM]}}
                \texttt{attempts = 1} \\
\texttt{target\_word = under} \\
\texttt{guess = slate} \\
\texttt{guess\_feedback = s<red> l<red> a<red> t<red> e<yellow>} \\
            }
        }
    }
    & & \\ \\

    \theutterance \stepcounter{utterance}  
    & & & \multicolumn{4}{p{0.6\linewidth}}{
        \cellcolor[rgb]{0.9,0.9,0.9}{
            \makecell[{{p{\linewidth}}}]{
                \texttt{\tiny{[P1$\langle$GM]}}
                \texttt{guess\_feedback: s<red> l<red> a<red> t<red> e<yellow>} \\
\\ 
\texttt{Provide your response only in this format.} \\
\texttt{explanation: details} \\
\texttt{guess: word} \\
            }
        }
    }
     \\ \\

    \theutterance \stepcounter{utterance}  
    & \multicolumn{4}{p{0.6\linewidth}}{
        \cellcolor[rgb]{0.561,0.737,0.561}{
            \makecell[{{p{\linewidth}}}]{
                \texttt{\tiny{[P1$\rangle$GM]}}
                \texttt{explanation: The feedback indicates that 's', 'l', 'a', 't' are not in the word, while 'e' is present but not in the last position. I will choose a word with 'e' in a different position and use other common letters like 'r', 'o', 'n', 'c' to test new possibilities.} \\
\texttt{guess: corey} \\
            }
        }
    }
    & & \\ \\

    \theutterance \stepcounter{utterance}  
    & & & \multicolumn{2}{p{0.3\linewidth}}{
        \cellcolor[rgb]{0.9,0.9,0.9}{
            \makecell[{{p{\linewidth}}}]{
                \texttt{\tiny{[GM$|$GM]}}
                \texttt{attempts = 2} \\
\texttt{target\_word = under} \\
\texttt{guess = corey} \\
\texttt{guess\_feedback = c<red> o<red> r<yellow> e<green> y<red>} \\
            }
        }
    }
    & & \\ \\

    \theutterance \stepcounter{utterance}  
    & & & \multicolumn{4}{p{0.6\linewidth}}{
        \cellcolor[rgb]{0.9,0.9,0.9}{
            \makecell[{{p{\linewidth}}}]{
                \texttt{\tiny{[P1$\langle$GM]}}
                \texttt{guess\_feedback: c<red> o<red> r<yellow> e<green> y<red>} \\
\\ 
\texttt{Provide your response only in this format.} \\
\texttt{explanation: details} \\
\texttt{guess: word} \\
            }
        }
    }
     \\ \\

    \theutterance \stepcounter{utterance}  
    & \multicolumn{4}{p{0.6\linewidth}}{
        \cellcolor[rgb]{0.561,0.737,0.561}{
            \makecell[{{p{\linewidth}}}]{
                \texttt{\tiny{[P1$\rangle$GM]}}
                \texttt{explanation: The feedback confirms 'e' is in the fourth position (green) and 'r' is in the word but not the third position (yellow). Letters 's', 'l', 'a', 't', 'c', 'o', 'y' are eliminated. I need a word with 'r' in position 1, 2, or 5, and 'e' in position 4, using remaining common letters like 'n', 'd', 'i', 'h'.} \\
\texttt{guess: dried} \\
            }
        }
    }
    & & \\ \\

    \theutterance \stepcounter{utterance}  
    & & & \multicolumn{2}{p{0.3\linewidth}}{
        \cellcolor[rgb]{0.9,0.9,0.9}{
            \makecell[{{p{\linewidth}}}]{
                \texttt{\tiny{[GM$|$GM]}}
                \texttt{attempts = 3} \\
\texttt{target\_word = under} \\
\texttt{guess = dried} \\
\texttt{guess\_feedback = d<yellow> r<yellow> i<red> e<green> d<red>} \\
            }
        }
    }
    & & \\ \\

    \theutterance \stepcounter{utterance}  
    & & & \multicolumn{4}{p{0.6\linewidth}}{
        \cellcolor[rgb]{0.9,0.9,0.9}{
            \makecell[{{p{\linewidth}}}]{
                \texttt{\tiny{[P1$\langle$GM]}}
                \texttt{guess\_feedback: d<yellow> r<yellow> i<red> e<green> d<red>} \\
\\ 
\texttt{Provide your response only in this format.} \\
\texttt{explanation: details} \\
\texttt{guess: word} \\
            }
        }
    }
     \\ \\

    \theutterance \stepcounter{utterance}  
    & \multicolumn{4}{p{0.6\linewidth}}{
        \cellcolor[rgb]{0.561,0.737,0.561}{
            \makecell[{{p{\linewidth}}}]{
                \texttt{\tiny{[P1$\rangle$GM]}}
                \texttt{explanation: The feedback confirms 'e' is in position 4. 'd' and 'r' are in the word but not in positions 1 and 2 respectively. 'i' is eliminated. Since 'd' was yellow in position 1 and red in position 5, it must be in position 3. 'r' was yellow in position 2, so it must be in position 1 or 5. With 'd' in position 3 and 'e' in position 4, the pattern is \_ \_ d e \_. 'r' cannot be in position 2 (already tested), so it must be in position 1 or 5. If 'r' is in position 1, we have r \_ d e \_. If 'r' is in position 5, we have \_ \_ d e r. Let's try 'r' in position 1 and 'd' in position 3, filling the rest with common letters like 'o' and 'n'. Wait, 'd' was yellow in pos 1, meaning it's in the word but not pos 1. It was red in pos 5, meaning it's not pos 5. So 'd' is in 2, 3, or 4. But 'e' is in 4. So 'd' is in 2 or 3. In 'dried', 'd' was at 1 (yellow) and 5 (red). So 'd' is in 2, 3, or 4. 'e' is 4. So 'd' is 2 or 3. 'r' was at 2 (yellow), so 'r' is 1, 3, 4, or 5. 'e' is 4. So 'r' is 1, 3, or 5. Let's try 'order'. 'o' is new, 'r' at 2 (was yellow in dried, so not 2), 'd' at 3, 'e' at 4, 'r' at 5. Wait, 'r' cannot be at 2. In 'order', 'r' is at 2 and 5. If 'r' is at 2, it would be yellow again? No, if the target has 'r' at 5, then 'r' at 2 is yellow. Let's try 'drone'. 'd' at 1 (yellow in dried, so not 1). 'r' at 2 (yellow in dried, so not 2). 'o' new. 'n' new. 'e' at 4 (green). This doesn't fit 'd' and 'r' constraints well. Let's re{-}evaluate. Target: \_ \_ \_ e \_. Known: 'e' at 4. 'd} \\
            }
        }
    }
    & & \\ \\

    \theutterance \stepcounter{utterance}  
    & & & \multicolumn{2}{p{0.3\linewidth}}{
        \cellcolor[rgb]{0.9,0.9,0.9}{
            \makecell[{{p{\linewidth}}}]{
                \texttt{\tiny{[GM$|$GM]}}
                \texttt{Error: The response should contain the 'guess:' keyword exactly once.} \\
            }
        }
    }
    & & \\ \\

    \theutterance \stepcounter{utterance}  
    & & & \multicolumn{2}{p{0.3\linewidth}}{
        \cellcolor[rgb]{0.9,0.9,0.9}{
            \makecell[{{p{\linewidth}}}]{
                \texttt{\tiny{[GM$|$GM]}}
                \texttt{game\_result = ABORT} \\
            }
        }
    }
    & & \\ \\

    \theutterance \stepcounter{utterance}  
    & & \multicolumn{4}{p{0.6\linewidth}}{
        \cellcolor[rgb]{0.95,0.95,0.95}{
            \makecell[{{p{\linewidth}}}]{
                \texttt{\tiny{[P1$\leftrightarrow$P]}}
                \texttt{Post{-}episode Reflexion (not sent to GM):} \\
\texttt{New plan: The previous attempt failed because I generated a long, rambling explanation that included internal monologue, self{-}correction, and incomplete sentences, violating the rule to provide a "concise explanation" and the implicit requirement to finish the response. I also failed to provide a valid "guess:" line with a complete five{-}letter word. Next time, I will strictly separate my reasoning from the output. I will first solve the logic internally to determine the exact word and a single{-}sentence justification. Then, I will output *only* the required format: a concise "explanation:" line summarizing the deduction (e.g., "Based on feedback, 'e' is at pos 4, 'd' at pos 3, and 'r' at pos 1, leading to the word 'drone'") followed immediately by "guess: {[}word{]}". I will ensure the word is exactly five letters, contains no eliminated letters, and fits all positional constraints before sending.} \\
            }
        }
    }
    & \\ \\

\end{supertabular}
}

\end{document}
